% !TeX program = lualatex
% !TeX root = predavanje12.tex
% ============================================================
% Zajednički preamble za sva predavanja RWA kolegija
% Uključiti s: % ============================================================
% Zajednički preamble za sva predavanja RWA kolegija
% Uključiti s: % ============================================================
% Zajednički preamble za sva predavanja RWA kolegija
% Uključiti s: \input{../common/preamble.tex}
% ============================================================

\documentclass[aspectratio=149,11pt]{beamer}
\usetheme{CambridgeUS}

\usepackage{fontspec}              % radi samo sa XeLaTeX/LuaLaTeX
\setsansfont{Latin Modern Sans}    % sans-serif font
\setmonofont{Latin Modern Mono}    % monospace font za kod
\usepackage[croatian]{babel}
\usepackage{csquotes}              % \enquote{} → hrvatski navodnici „..."
\usepackage{lmodern}
\usepackage{booktabs}
\usepackage{listings}
\usepackage{xcolor}
\usepackage{hyperref}

% ---------- Autor i logo ----------
\author{mag.ing. Josip Torić (SIT UNIZD)}
\titlegraphic{\includegraphics[height=2.5cm]{../common/SIT_logo_tekst1.pdf}}

% ---------- Listings ----------
\definecolor{codebg}{RGB}{248,248,248}
\definecolor{codeframeclr}{RGB}{220,220,220}
\lstdefinelanguage{bash}{
  morekeywords={sudo,apt,brew,pip,pipx,uv,winget,python,git,flask,source,deactivate,code,gunicorn,docker,compose,heroku,render,railway,node,npm,pnpm,yarn,uvicorn,alembic,pytest,ruff,black,precommit},
  sensitive=false,
  morecomment=[l]{\#}
}
\lstdefinelanguage{python}{
  morekeywords={from,import,def,class,return,if,elif,else,for,in,while,try,except,raise,with,as,not,and,or,is,lambda,yield,True,False,None,await,async},
  sensitive=true,
  morecomment=[l]{\#},
  morestring=[b]',
  morestring=[b]"
}
\lstdefinelanguage{sql}{
  morekeywords={SELECT,FROM,WHERE,INSERT,INTO,VALUES,UPDATE,SET,DELETE,CREATE,TABLE,PRIMARY,KEY,FOREIGN,REFERENCES,UNIQUE,NOT,NULL,INDEX,ON,ALTER,DROP,CONSTRAINT},
  sensitive=false,
  morecomment=[l]{--}
}
\lstdefinelanguage{json}{
  morecomment=[l]{//},
  morestring=[b]",
  showstringspaces=false
}
\lstdefinelanguage{typescript}{
  morekeywords={import,export,from,const,let,var,function,return,if,else,for,of,in,while,switch,case,break,async,await,try,catch,finally,throw,new,class,this,extends,implements,interface,type,enum,readonly,public,private,protected,static,get,set,true,false,null,undefined,void,string,number,boolean,any,unknown,never,typeof,keyof,as},
  sensitive=true,
  morecomment=[l]{//},
  morecomment=[s]{/*}{*/},
  morestring=[b]',
  morestring=[b]",
  morestring=[b]`
}
\lstalias{javascript}{typescript}
\lstset{
  backgroundcolor=\color{codebg},
  basicstyle=\ttfamily\footnotesize,
  frame=single,
  rulecolor=\color{codeframeclr},
  frameround=tttt,
  breaklines=true,
  tabsize=2,
  showstringspaces=false,
  upquote=true,
  columns=fullflexible
}

% ---------- Helper macros ----------
\newcommand{\codefile}[1]{\texttt{#1}}
\newcommand{\kbd}[1]{\texttt{#1}}

% ---------- Code frame environment ----------
% Automatski dodaje fragile i shrink=5, koristi se umjesto \begin{frame}
% Primjer: \begin{codeframe}{Naslov} ... \end{codeframe}
% Primjer s opcijama: \begin{codeframe}[shrink=15]{Naslov} ... \end{codeframe}
\newenvironment{codeframe}[2][]
  {\begin{frame}[fragile,shrink=5,#1]{#2}}
  {\end{frame}}

% ============================================================

\documentclass[aspectratio=149,11pt]{beamer}
\usetheme{CambridgeUS}

\usepackage{fontspec}              % radi samo sa XeLaTeX/LuaLaTeX
\setsansfont{Latin Modern Sans}    % sans-serif font
\setmonofont{Latin Modern Mono}    % monospace font za kod
\usepackage[croatian]{babel}
\usepackage{csquotes}              % \enquote{} → hrvatski navodnici „..."
\usepackage{lmodern}
\usepackage{booktabs}
\usepackage{listings}
\usepackage{xcolor}
\usepackage{hyperref}

% ---------- Autor i logo ----------
\author{mag.ing. Josip Torić (SIT UNIZD)}
\titlegraphic{\includegraphics[height=2.5cm]{../common/SIT_logo_tekst1.pdf}}

% ---------- Listings ----------
\definecolor{codebg}{RGB}{248,248,248}
\definecolor{codeframeclr}{RGB}{220,220,220}
\lstdefinelanguage{bash}{
  morekeywords={sudo,apt,brew,pip,pipx,uv,winget,python,git,flask,source,deactivate,code,gunicorn,docker,compose,heroku,render,railway,node,npm,pnpm,yarn,uvicorn,alembic,pytest,ruff,black,precommit},
  sensitive=false,
  morecomment=[l]{\#}
}
\lstdefinelanguage{python}{
  morekeywords={from,import,def,class,return,if,elif,else,for,in,while,try,except,raise,with,as,not,and,or,is,lambda,yield,True,False,None,await,async},
  sensitive=true,
  morecomment=[l]{\#},
  morestring=[b]',
  morestring=[b]"
}
\lstdefinelanguage{sql}{
  morekeywords={SELECT,FROM,WHERE,INSERT,INTO,VALUES,UPDATE,SET,DELETE,CREATE,TABLE,PRIMARY,KEY,FOREIGN,REFERENCES,UNIQUE,NOT,NULL,INDEX,ON,ALTER,DROP,CONSTRAINT},
  sensitive=false,
  morecomment=[l]{--}
}
\lstdefinelanguage{json}{
  morecomment=[l]{//},
  morestring=[b]",
  showstringspaces=false
}
\lstdefinelanguage{typescript}{
  morekeywords={import,export,from,const,let,var,function,return,if,else,for,of,in,while,switch,case,break,async,await,try,catch,finally,throw,new,class,this,extends,implements,interface,type,enum,readonly,public,private,protected,static,get,set,true,false,null,undefined,void,string,number,boolean,any,unknown,never,typeof,keyof,as},
  sensitive=true,
  morecomment=[l]{//},
  morecomment=[s]{/*}{*/},
  morestring=[b]',
  morestring=[b]",
  morestring=[b]`
}
\lstalias{javascript}{typescript}
\lstset{
  backgroundcolor=\color{codebg},
  basicstyle=\ttfamily\footnotesize,
  frame=single,
  rulecolor=\color{codeframeclr},
  frameround=tttt,
  breaklines=true,
  tabsize=2,
  showstringspaces=false,
  upquote=true,
  columns=fullflexible
}

% ---------- Helper macros ----------
\newcommand{\codefile}[1]{\texttt{#1}}
\newcommand{\kbd}[1]{\texttt{#1}}

% ---------- Code frame environment ----------
% Automatski dodaje fragile i shrink=5, koristi se umjesto \begin{frame}
% Primjer: \begin{codeframe}{Naslov} ... \end{codeframe}
% Primjer s opcijama: \begin{codeframe}[shrink=15]{Naslov} ... \end{codeframe}
\newenvironment{codeframe}[2][]
  {\begin{frame}[fragile,shrink=5,#1]{#2}}
  {\end{frame}}

% ============================================================

\documentclass[aspectratio=149,11pt]{beamer}
\usetheme{CambridgeUS}

\usepackage{fontspec}              % radi samo sa XeLaTeX/LuaLaTeX
\setsansfont{Latin Modern Sans}    % sans-serif font
\setmonofont{Latin Modern Mono}    % monospace font za kod
\usepackage[croatian]{babel}
\usepackage{csquotes}              % \enquote{} → hrvatski navodnici „..."
\usepackage{lmodern}
\usepackage{booktabs}
\usepackage{listings}
\usepackage{xcolor}
\usepackage{hyperref}

% ---------- Autor i logo ----------
\author{mag.ing. Josip Torić (SIT UNIZD)}
\titlegraphic{\includegraphics[height=2.5cm]{../common/SIT_logo_tekst1.pdf}}

% ---------- Listings ----------
\definecolor{codebg}{RGB}{248,248,248}
\definecolor{codeframeclr}{RGB}{220,220,220}
\lstdefinelanguage{bash}{
  morekeywords={sudo,apt,brew,pip,pipx,uv,winget,python,git,flask,source,deactivate,code,gunicorn,docker,compose,heroku,render,railway,node,npm,pnpm,yarn,uvicorn,alembic,pytest,ruff,black,precommit},
  sensitive=false,
  morecomment=[l]{\#}
}
\lstdefinelanguage{python}{
  morekeywords={from,import,def,class,return,if,elif,else,for,in,while,try,except,raise,with,as,not,and,or,is,lambda,yield,True,False,None,await,async},
  sensitive=true,
  morecomment=[l]{\#},
  morestring=[b]',
  morestring=[b]"
}
\lstdefinelanguage{sql}{
  morekeywords={SELECT,FROM,WHERE,INSERT,INTO,VALUES,UPDATE,SET,DELETE,CREATE,TABLE,PRIMARY,KEY,FOREIGN,REFERENCES,UNIQUE,NOT,NULL,INDEX,ON,ALTER,DROP,CONSTRAINT},
  sensitive=false,
  morecomment=[l]{--}
}
\lstdefinelanguage{json}{
  morecomment=[l]{//},
  morestring=[b]",
  showstringspaces=false
}
\lstdefinelanguage{typescript}{
  morekeywords={import,export,from,const,let,var,function,return,if,else,for,of,in,while,switch,case,break,async,await,try,catch,finally,throw,new,class,this,extends,implements,interface,type,enum,readonly,public,private,protected,static,get,set,true,false,null,undefined,void,string,number,boolean,any,unknown,never,typeof,keyof,as},
  sensitive=true,
  morecomment=[l]{//},
  morecomment=[s]{/*}{*/},
  morestring=[b]',
  morestring=[b]",
  morestring=[b]`
}
\lstalias{javascript}{typescript}
\lstset{
  backgroundcolor=\color{codebg},
  basicstyle=\ttfamily\footnotesize,
  frame=single,
  rulecolor=\color{codeframeclr},
  frameround=tttt,
  breaklines=true,
  tabsize=2,
  showstringspaces=false,
  upquote=true,
  columns=fullflexible
}

% ---------- Helper macros ----------
\newcommand{\codefile}[1]{\texttt{#1}}
\newcommand{\kbd}[1]{\texttt{#1}}

% ---------- Code frame environment ----------
% Automatski dodaje fragile i shrink=5, koristi se umjesto \begin{frame}
% Primjer: \begin{codeframe}{Naslov} ... \end{codeframe}
% Primjer s opcijama: \begin{codeframe}[shrink=15]{Naslov} ... \end{codeframe}
\newenvironment{codeframe}[2][]
  {\begin{frame}[fragile,shrink=5,#1]{#2}}
  {\end{frame}}


\title[Programiranje za web]{Programiranje za web - Predavanje 12 \\ Production, deploy i završni polish}
\date{11.01.2027.}

\begin{document}

\frame{\titlepage}

% ------------------------------------------------------------
\begin{frame}{Plan predavanja}
\begin{enumerate}
  \item Razlika dev vs production okruženje
  \item Environment varijable i konfiguracija
  \item CORS: što je i zašto treba
  \item Deploy backenda na Railway
  \item Deploy frontenda (statički build)
  \item CI/CD pipeline
  \item Priprema za 2. obranu
  \item Commitovi za ovo predavanje
\end{enumerate}

\vspace{0.6em}
Na kraju predavanja: aplikacija je dostupna na javnom URL-u, s CI pipelineom koji provjerava kod.
\end{frame}

% ------------------------------------------------------------
\begin{frame}{Ishodi učenja}
Nakon ovog predavanja student može:
\begin{itemize}
  \item objasniti razliku između dev i production konfiguracije
  \item konfigurirati CORS na backendu za komunikaciju s frontendom
  \item deployati FastAPI backend na Railway ili Render
  \item konfigurirati Vite build za produkciju s environment varijablama
  \item postaviti osnovni CI pipeline za lint, testove i build
\end{itemize}
\end{frame}

% ------------------------------------------------------------
\begin{frame}{Dev vs Production}
\begin{tabular}{@{}lp{5cm}p{5cm}@{}}
  \toprule
  & \textbf{Development} & \textbf{Production} \\
  \midrule
  Backend & \texttt{uvicorn --reload}, localhost:8000 & Gunicorn + Uvicorn workers, javni URL \\
  Frontend & Vite dev server, localhost:5173 & Statički build (\texttt{dist/}) \\
  Baza & Docker lokalno & Managed PostgreSQL \\
  Debug & Swagger UI, detalji grešaka & Bez Swaggera, generičke greške \\
  Env & \texttt{.env} file & Platform env vars \\
  \bottomrule
\end{tabular}

\vspace{0.6em}
Nikad ne deployajte s \texttt{--reload} ili s \texttt{DEBUG=True}. Production konfiguracija skriva detalje grešaka od korisnika.
\end{frame}

% ------------------------------------------------------------
\begin{frame}{Production mindset}
Razvoj i produkcija zahtijevaju različit način razmišljanja:

\vspace{0.4em}
\begin{tabular}{@{}lp{5cm}p{5cm}@{}}
  \toprule
  & \textbf{Development} & \textbf{Production} \\
  \midrule
  Greške & Detaljan stack trace & Generička poruka korisniku, detalji u logovima \\
  Padanje & Prihvatljivo - restart & Mora se oporaviti automatski \\
  Podaci & Seed podaci, testni & Pravi korisnici, pravi podaci \\
  Sigurnost & Relaksirano & Sve zaključano \\
  \bottomrule
\end{tabular}

\vspace{0.6em}
Princip: \textbf{fail gracefully}. Umjesto bijelog ekrana, korisnik vidi \enquote{Nešto je pošlo po krivu, pokušajte kasnije}. Detalji greške idu u logove, ne korisniku.
\end{frame}

% ------------------------------------------------------------
\begin{frame}{Security checklist za deploy}
\begin{enumerate}
  \item \textbf{HTTPS svugdje} - nikad HTTP u produkciji (Railway/Render to rade automatski)
  \item \textbf{Env vars za tajne} - JWT\_SECRET, DATABASE\_URL nikad u kodu
  \item \textbf{Jaki JWT\_SECRET} - minimalno 32 znaka, random generirani
  \item \textbf{CORS ograničen} - samo vaš frontend origin, ne \texttt{*}
  \item \textbf{Swagger isključen} - \texttt{docs\_url=None} u produkciji
  \item \textbf{Dependency audit} - \texttt{pip audit} i \texttt{npm audit} za poznate ranjivosti
\end{enumerate}
\end{frame}

% ------------------------------------------------------------
\begin{frame}{12-Factor App: relevantni principi}
\begin{tabular}{@{}lp{8cm}@{}}
  \toprule
  \textbf{Princip} & \textbf{Što znači za nas} \\
  \midrule
  Config u envu & Sve konfigurabilno kroz env varijable, ne hardkodirano \\
  Stateless procesi & Server ne čuva stanje između requestova (session je u tokenu) \\
  Port binding & Aplikacija sluša na portu koji joj je zadan (\texttt{\$PORT}) \\
  Dev/prod paritet & Lokalni dev je što sličniji produkciji \\
  Logovi kao stream & \texttt{print}/\texttt{logging} na stdout, platforma ih skuplja \\
  \bottomrule
\end{tabular}

\vspace{0.6em}
Ovi principi nisu akademski - Railway i Render ih očekuju. Aplikacija koja ih slijedi deploya se bez problema, aplikacija koja ih ne slijedi zahtijeva workarounde.
\end{frame}

% ------------------------------------------------------------
\begin{frame}[fragile,shrink=5]{Environment varijable: backend}
\begin{lstlisting}[language=python]
# app/core/config.py
from pydantic_settings import (
    BaseSettings, SettingsConfigDict
)

class Settings(BaseSettings):
    model_config = SettingsConfigDict(env_file=".env")

    DATABASE_URL: str
    JWT_SECRET: str
    CORS_ORIGINS: str = "http://localhost:5173"
    DEBUG: bool = False

settings = Settings()
\end{lstlisting}

\vspace{0.4em}
Pydantic \texttt{BaseSettings} automatski čita iz env varijabli. U developmentu koristi \texttt{.env} file, u produkciji čita iz platform env vars.

\vspace{0.4em}
\texttt{.env} nikad ne smije biti committan u git - stavite ga u \texttt{.gitignore}. U repozitoriju držite \texttt{.env.example} s praznim vrijednostima.
\end{frame}

% ------------------------------------------------------------
\begin{frame}[fragile,shrink=5]{Environment varijable: frontend (Vite)}
\begin{lstlisting}[language=bash]
# .env.development
VITE_API_URL=http://localhost:8000

# .env.production
VITE_API_URL=https://moj-api.railway.app
\end{lstlisting}

\begin{lstlisting}[language=typescript]
// src/api/client.ts
const api = axios.create({
  baseURL: import.meta.env.VITE_API_URL,
})
\end{lstlisting}

\vspace{0.4em}
Vite automatski učitava \texttt{.env.development} u dev modu i \texttt{.env.production} za build. Sve varijable moraju počinjati s \texttt{VITE\_} da budu dostupne u kodu.

\vspace{0.4em}
Ove varijable se ugrađuju u build (\texttt{npm run build}) - nisu runtime, nego build-time.
\end{frame}

% ------------------------------------------------------------
\begin{frame}{CORS: Cross-Origin Resource Sharing}
Browser blokira requestove s jedne domene na drugu (Same-Origin Policy). CORS je mehanizam koji dopušta iznimke.

\vspace{0.4em}
\begin{tabular}{@{}lll@{}}
  \toprule
  \textbf{Scenario} & \textbf{Origin} & \textbf{Blocked?} \\
  \midrule
  Dev frontend $\rightarrow$ dev backend & :5173 $\rightarrow$ :8000 & Da (različiti portovi) \\
  Prod frontend $\rightarrow$ prod backend & app.com $\rightarrow$ api.app.com & Da (različita subdomena) \\
  Postman $\rightarrow$ backend & nema origin & Ne (nije browser) \\
  \bottomrule
\end{tabular}

\vspace{0.6em}
Backend mora eksplicitno dopustiti origine s kojih prima requestove. Bez CORS konfiguracije, frontend neće moći komunicirati s API-jem iz browsera.
\end{frame}

% ------------------------------------------------------------
\begin{frame}[fragile,shrink=5]{CORS konfiguracija u FastAPI}
\begin{lstlisting}[language=python]
# app/main.py
from fastapi.middleware.cors import CORSMiddleware
from app.core.config import settings

app.add_middleware(
    CORSMiddleware,
    allow_origins=settings.CORS_ORIGINS.split(","),
    allow_credentials=True,
    allow_methods=["*"],
    allow_headers=["*"],
)
\end{lstlisting}

\vspace{0.4em}
\begin{tabular}{@{}ll@{}}
  \toprule
  \textbf{Parametar} & \textbf{Značenje} \\
  \midrule
  \texttt{allow\_origins} & Koji origini smiju slati requestove \\
  \texttt{allow\_credentials} & Smiju li se slati cookieji \\
  \texttt{allow\_methods} & Koje HTTP metode (GET, POST...) \\
  \texttt{allow\_headers} & Koji headeri (Authorization...) \\
  \bottomrule
\end{tabular}

\vspace{0.4em}
Nikad ne koristite \texttt{allow\_origins=["*"]} u produkciji - to dopušta bilo kojoj stranici da šalje requestove na vaš API.
\end{frame}

% ------------------------------------------------------------
\begin{frame}{Preflight request: kad i zašto}
Browser šalje \textbf{OPTIONS} request prije stvarnog requesta u određenim situacijama:

\vspace{0.4em}
\begin{tabular}{@{}lp{4cm}p{5cm}@{}}
  \toprule
  & \textbf{Simple request} & \textbf{Preflighted request} \\
  \midrule
  Metode & GET, POST, HEAD & PUT, PATCH, DELETE \\
  Headeri & Samo standardni & Custom headeri (Authorization) \\
  Content-Type & form-data, text/plain & application/json \\
  Preflight & Ne & Da - OPTIONS prvo \\
  \bottomrule
\end{tabular}

\vspace{0.6em}
Naši API pozivi šalju \texttt{Authorization} header i \texttt{Content-Type: application/json} - svaki request je preflighted. Browser prvo pita \enquote{smijem li?} (OPTIONS), backend odgovori s CORS headerima, tek onda ide stvarni request.
\end{frame}

% ------------------------------------------------------------
\begin{frame}{Content Security Policy}
CSP je HTTP header koji kaže browseru odakle smije učitavati resurse:

\vspace{0.4em}
\begin{itemize}
  \item \textbf{Bez CSP-a}: maliciozna skripta može učitati JavaScript s bilo kojeg URL-a
  \item \textbf{S CSP-om}: browser odbija sve resurse koji nisu s odobrenih origina
\end{itemize}

\vspace{0.6em}
\begin{tabular}{@{}lp{8cm}@{}}
  \toprule
  \textbf{Direktiva} & \textbf{Kontrolira} \\
  \midrule
  \texttt{default-src 'self'} & Sve resurse s istog origina \\
  \texttt{script-src 'self'} & JavaScript samo s našeg servera \\
  \texttt{connect-src} & API pozivi - ovdje dodajemo backend URL \\
  \texttt{img-src} & Odakle smiju dolaziti slike \\
  \bottomrule
\end{tabular}

\vspace{0.4em}
CSP je dodatni sloj zaštite od XSS-a. Čak i ako napadač ubaci skriptu, browser je neće izvršiti ako origin nije odobren.
\end{frame}

% ------------------------------------------------------------
\begin{frame}{Deploy opcije}
\begin{tabular}{@{}llll@{}}
  \toprule
  \textbf{Platforma} & \textbf{Backend} & \textbf{Frontend} & \textbf{Baza} \\
  \midrule
  Railway & Python buildpack & Static site & Railway PostgreSQL \\
  Render & Python, Docker & Static site & Render PostgreSQL \\
  Fly.io & Docker & CDN & Fly PostgreSQL \\
  \bottomrule
\end{tabular}

\vspace{0.6em}
Railway i Render imaju besplatne tierove dovoljne za studentski projekt. Odaberite jednu platformu i držite se nje.
\end{frame}

% ------------------------------------------------------------
\begin{frame}[fragile,shrink=5]{Deploy backenda na Railway}
\begin{lstlisting}[language=bash]
# 1. Dodajte Procfile u api/
# Procfile
web: uvicorn app.main:app --host 0.0.0.0 --port $PORT

# 2. requirements.txt: rucno odrzavana lista
#    direktnih ovisnosti (fastapi, uvicorn,
#    sqlalchemy, asyncpg, alembic, pyjwt, bcrypt...)
#    Ne generirajte ga s "uv pip freeze" - to zaledi
#    i tranzitivne ovisnosti i tesko se odrzava.

# 3. Na Railway:
#    - New Project > Deploy from GitHub repo
#    - Root directory: api
#    - Add PostgreSQL plugin
#    - Set env vars: DATABASE_URL, JWT_SECRET, CORS_ORIGINS
\end{lstlisting}

\vspace{0.4em}
Railway automatski detektira Python projekt, instalira dependencije i pokreće \texttt{Procfile} naredbu. \texttt{\$PORT} je env varijabla koju Railway postavlja.
\end{frame}

% ------------------------------------------------------------
\begin{frame}[fragile,shrink=5]{Deploy frontenda}
\begin{lstlisting}[language=bash]
# 1. Build za produkciju
cd web
npm run build     # generira dist/ folder

# 2. dist/ sadrzi staticke fileove:
#    index.html, assets/main-abc123.js, assets/style-def456.css

# 3. Na Railway/Render:
#    - New Static Site > Deploy from GitHub
#    - Build command: npm run build
#    - Publish directory: dist
#    - Set VITE_API_URL env var
\end{lstlisting}

\vspace{0.4em}
SPA routing zahtijeva konfiguraciju: sve rute moraju servirati \texttt{index.html}. Na Railway/Render postavite \textbf{fallback} na \texttt{index.html} za 404 stranice.
\end{frame}

% ------------------------------------------------------------
\begin{frame}{CDN i caching strategija}
Statički asseti (JS, CSS, slike) se serviraju s CDN-a za brži pristup:

\vspace{0.4em}
\begin{itemize}
  \item Vite generira hashirane nazive: \texttt{main-abc123.js}
  \item Hash se mijenja samo kad se kod promijeni
  \item Server može cachirati hashirane fileove zauvijek (\texttt{Cache-Control: immutable})
  \item Novi deploy = novi hash = browser automatski dohvati novu verziju
\end{itemize}

\vspace{0.6em}
\texttt{index.html} se \textbf{nikad ne cachira} (ili s kratkim max-age) jer sadrži linkove na hashirane assete. Kad se deploy promijeni, novi index.html referencira nove fileove.
\end{frame}

% ------------------------------------------------------------
\begin{frame}{Build optimizacije}
\texttt{npm run build} generira optimizirani bundle za produkciju:

\vspace{0.4em}
\begin{tabular}{@{}lp{8cm}@{}}
  \toprule
  \textbf{Optimizacija} & \textbf{Što radi} \\
  \midrule
  Tree shaking & Uklanja nekorišteni kod iz bundlea \\
  Code splitting & Lazy-loaded rute postaju zasebni chunkovi \\
  Minification & Skraćuje nazive varijabli, uklanja whitespace \\
  CSS extraction & CSS u zasebnom fileu (paralelno učitavanje) \\
  \bottomrule
\end{tabular}

\vspace{0.6em}
Rezultat: umjesto jednog velikog JS filea, dobijemo glavni chunk + chunk po ruti. Korisnik preuzima samo kod za stranicu koju gleda.
\end{frame}

% ------------------------------------------------------------
\begin{frame}{Performance budžet}
Koliko JavaScripta je previše?

\vspace{0.4em}
\begin{tabular}{@{}ll@{}}
  \toprule
  \textbf{Bundle veličina} & \textbf{Procjena} \\
  \midrule
  < 100 KB (gzipped) & Odlično \\
  100-250 KB & Prihvatljivo za SPA \\
  250-500 KB & Treba optimizirati \\
  > 500 KB & Problem - korisnici na sporom internetu čekaju \\
  \bottomrule
\end{tabular}

\vspace{0.6em}
Alati za analizu: \texttt{npx vite-bundle-visualizer} pokazuje koji paketi zauzimaju koliko prostora. Često su krivci: moment.js, lodash (cijeli), UI biblioteke. Vue + Vue Router + Pinia + Axios zajedno su oko 50 KB gzipped - razumno.
\end{frame}

% ------------------------------------------------------------
\begin{frame}{SPA fallback: zašto je potreban}
Kad korisnik direktno upiše \texttt{app.com/clubs/5/lifters}:

\vspace{0.4em}
\begin{enumerate}
  \item Browser šalje GET \texttt{/clubs/5/lifters} web serveru
  \item Web server traži file \texttt{clubs/5/lifters/index.html} - ne postoji
  \item Bez fallbacka: \textbf{404 Not Found}
  \item S fallbackom: servira \texttt{index.html} $\rightarrow$ Vue Router preuzima $\rightarrow$ renderira \texttt{LiftersView}
\end{enumerate}

\vspace{0.6em}
Ovo je razlika između MPA (server zna rute) i SPA (klijent zna rute). Server mora sve nepoznate rute preusmjeriti na \texttt{index.html}.
\end{frame}

% ------------------------------------------------------------
\begin{frame}[fragile,shrink=5]{Migracije na produkcijskoj bazi}
\begin{lstlisting}[language=bash]
# Na Railway: dodajte u start komandu
# Procfile
release: cd api && alembic upgrade head
web: uvicorn app.main:app --host 0.0.0.0 --port $PORT
\end{lstlisting}

\vspace{0.4em}
\texttt{release} komanda se izvršava \textbf{prije} pokretanja aplikacije, pri svakom deployu. Alembic automatski primijeni sve pending migracije.

\vspace{0.6em}
Nikad ručno ne mijenjajte production bazu. Sve promjene idu kroz migracije:
\begin{enumerate}
  \item Lokalno: \texttt{alembic revision --autogenerate -m "opis"}
  \item Testirajte lokalno: \texttt{alembic upgrade head}
  \item Commitajte migraciju
  \item Deploy - \texttt{release} komanda automatski pokrene migraciju
\end{enumerate}
\end{frame}

% ------------------------------------------------------------
\begin{frame}[fragile,shrink=5]{CI pipeline: GitHub Actions}
\begin{lstlisting}[language=bash]
# .github/workflows/ci.yml
name: CI
on: [push, pull_request]

jobs:
  backend:
    runs-on: ubuntu-latest
    steps:
      - uses: actions/checkout@v4
      - uses: actions/setup-python@v5
        with:
          python-version: "3.12"
      - uses: astral-sh/setup-uv@v5
      - run: uv pip install --system -r api/requirements.txt
      - run: cd api && ruff check . && ruff format --check .
      - run: cd api && pytest

  frontend:
    runs-on: ubuntu-latest
    steps:
      - uses: actions/checkout@v4
      - uses: actions/setup-node@v4
        with:
          node-version: "22"
      - run: cd web && npm ci
      - run: cd web && npm run lint
      - run: cd web && npm run type-check
      - run: cd web && npm run test -- --run
      - run: cd web && npm run build
\end{lstlisting}
\end{frame}

% ------------------------------------------------------------
\begin{frame}{CI/CD: što i zašto}
\begin{tabular}{@{}lp{10cm}@{}}
  \toprule
  \textbf{Korak} & \textbf{Što provjerava} \\
  \midrule
  Lint (ruff/eslint) & Formatiranje, stilske greške, potencijalni bugovi \\
  Testovi & pytest i Vitest iz Predavanja 11 - logika radi, nema regresija \\
  Build (npm run build) & Frontend se kompajlira bez grešaka \\
  \bottomrule
\end{tabular}

\vspace{0.6em}
Pipeline se pokreće na svaki push i PR. Ako bilo koji korak padne, commit je označen crvenim X-om. Ovako broken kod ne dolazi u produkciju.

\vspace{0.4em}
Za ovaj kolegij CI nije obavezan, ali je bonus. Ako ga postavite, pokazuje zrelost projekta na obrani.
\end{frame}

% ------------------------------------------------------------
\begin{frame}{CD: Continuous Deployment}
CI provjerava kvalitetu koda. CD automatski deploya kod koji prođe provjere:

\vspace{0.4em}
\begin{enumerate}
  \item Push na \texttt{main} branch
  \item CI pokrene lint + testove + build
  \item Ako sve prođe $\rightarrow$ automatski deploy na produkciju
  \item Ako padne $\rightarrow$ deploy se ne dogodi, developer dobije obavijest
\end{enumerate}

\vspace{0.6em}
Railway i Render automatski deployaju svaki push na main - to je CD bez ikakve konfiguracije. Za složenije sustave koristimo staging environment (test server) prije produkcije.

\vspace{0.4em}
Rollback strategija: ako novi deploy pokvari aplikaciju, revertamo commit i pushamo - automatski se deploya prethodna verzija.
\end{frame}

% ------------------------------------------------------------
\begin{frame}{Monitoring u produkciji}
Aplikacija u produkciji mora biti nadgledana:

\vspace{0.4em}
\begin{tabular}{@{}lp{4cm}p{5cm}@{}}
  \toprule
  \textbf{Što} & \textbf{Alat} & \textbf{Zašto} \\
  \midrule
  Error tracking & Sentry & Hvata exceptions i šalje alert \\
  Uptime monitoring & UptimeRobot & Provjera je li server živ (health check) \\
  Logovi & Railway logs & Debug problema u produkciji \\
  Performance & Lighthouse & Brzina učitavanja frontenda \\
  \bottomrule
\end{tabular}

\vspace{0.6em}
Bez monitoringa ne znate da aplikacija ne radi dok vam se student ne javi. Health check endpoint (\texttt{GET /health}) je minimum - monitoring alat ga pinguje svake minute.
\end{frame}

% ------------------------------------------------------------
\begin{frame}{Logging strategija}
\begin{itemize}
  \item \textbf{Strukturirani logovi}: JSON format s timestamp, level, message, context
  \item \textbf{Log levels}: DEBUG (dev), INFO (normalne operacije), WARNING (sumnjivo), ERROR (greška)
  \item \textbf{Što logirati}: login pokušaji, CRUD akcije, greške, spori upiti
  \item \textbf{Što ne logirati}: lozinke, tokene, osobne podatke (GDPR)
\end{itemize}

\vspace{0.6em}
U produkciji logovi idu na stdout - platforma (Railway) ih skuplja i prikazuje u dashboardu. Ne pišite u fileove na serveru - kontejneri su efemerni.
\end{frame}

% ------------------------------------------------------------
\begin{frame}{Checklist prije deploya}
\begin{enumerate}
  \item \texttt{.env} je u \texttt{.gitignore} (nikad committan)
  \item \texttt{.env.example} postoji s praznim vrijednostima
  \item CORS origins postavljeni na frontend URL (ne \texttt{*})
  \item Swagger UI isključen u produkciji (\texttt{docs\_url=None} kad \texttt{DEBUG=False})
  \item \texttt{JWT\_SECRET} je dugačak random string (ne \texttt{secret123})
  \item Frontend build radi: \texttt{npm run build} bez grešaka
  \item \texttt{VITE\_API\_URL} postavljen na production backend URL
  \item Migracije se automatski pokreću pri deployu
  \item Health check endpoint radi: \texttt{GET /health}
\end{enumerate}
\end{frame}

% ------------------------------------------------------------
\begin{frame}{Database backup i recovery}
Podaci u produkciji su najvredniji dio aplikacije:

\vspace{0.4em}
\begin{itemize}
  \item \textbf{Automatizirani backupi}: Railway/Render rade dnevne backupe managed PostgreSQL-a
  \item \textbf{Point-in-time recovery}: mogućnost vraćanja baze na bilo koji trenutak u zadnjih X dana
  \item \textbf{Testiranje restora}: backup koji nikad niste testirali nije backup
\end{itemize}

\vspace{0.6em}
Scenarij: pogrešna migracija obriše stupac s podacima. Bez backupa - izgubljeni podaci. S backupom - restore na stanje prije migracije. Zato nikad ne radite \texttt{DROP COLUMN} bez prethodne provjere i backupa.
\end{frame}

% ------------------------------------------------------------
\begin{frame}[shrink=10]{Što dalje nakon kolegija}
Tehnologije koje nadograđuju ono što ste naučili:

\vspace{0.4em}
\begin{tabular}{@{}lp{8cm}@{}}
  \toprule
  \textbf{Tehnologija} & \textbf{Zašto} \\
  \midrule
  Nuxt & Server-side rendering za Vue - SEO, brži first paint \\
  Docker & Kontejnerizacija: isti environment lokalno i na serveru \\
  WebSockets & Real-time komunikacija bez pollinga (chat, live score) \\
  PWA & Progressive Web App - offline support, instalacija na mobitel \\
  GraphQL / tRPC & Type-safe API komunikacija bez pisanja duplog koda \\
  TanStack Query & Server-state s cachingom i retryjem - zamjenjuje ručni loading/error kod iz P9 \\
  \bottomrule
\end{tabular}

\vspace{0.6em}
TypeScript i testiranje (Vitest, Playwright) su dio ovog kolegija. Ono što ste naučili je solidan temelj - svaka od ovih tehnologija nadograđuje iste principe.
\end{frame}

% ------------------------------------------------------------
\begin{frame}{Commitovi za Predavanje 12}
\begin{itemize}
  \item \texttt{chore: configure production env and CORS settings for deployment}
  \item \texttt{feat(web): add build-time env config and runtime API base handling}
  \item \texttt{ci: add pipeline for lint/test/build (api + web)}
  \item \texttt{docs: add deploy guide (Railway) and production checklist}
  \item \texttt{refactor: final cleanup, remove dead code, tighten types}
\end{itemize}

\vspace{0.4em}
\textbf{Na satu:} deployamo aplikaciju na Railway, gledamo logove, testiramo javni URL.
\end{frame}

% ------------------------------------------------------------
\begin{frame}{Priprema za 2. obranu (18.01.)}
\begin{enumerate}
  \item Aplikacija je deployana i dostupna na javnom URL-u
  \item Demonstrirate: login, CRUD, pravila, error stanja
  \item Pokažete Network tab u DevToolsima - requestovi i responsovi
  \item Budite spremni objasniti:
  \begin{itemize}
    \item Kako frontend hendla 401? Gdje je interceptor?
    \item Što se dogodi kad korisnik ručno upiše admin URL?
    \item Gdje držite stanje? Zašto store, a ne lokalni ref?
    \item Kako ste konfigurirali CORS i zašto?
    \item Kako radi deploy pipeline?
  \end{itemize}
  \item Mini promjena uživo: dodaj polje u formu, prikaži novi podatak
\end{enumerate}
\end{frame}

% ------------------------------------------------------------
\begin{frame}{Rekapitulacija: cijeli stack}
\begin{columns}[T]
\column{0.48\textwidth}
\textbf{Backend (P1-P6)}
\begin{enumerate}
  \item Skeleton, config, health
  \item Modeli, ORM, migracije
  \item Autentikacija (JWT)
  \item Autorizacija (role, ownership)
  \item CRUD, validacija, paginacija
  \item Poslovna pravila, workflow
\end{enumerate}

\column{0.48\textwidth}
\textbf{Frontend (P7-P12)}
\begin{enumerate}
  \setcounter{enumi}{6}
  \item SPA, router, layout
  \item API komunikacija, auth flow
  \item State management, forme
  \item Admin tablice, filteri, UI
  \item Testiranje: Vitest, Playwright
  \item Deploy, CORS, CI/CD
\end{enumerate}
\end{columns}

\vspace{0.4em}
12 predavanja, jedan cjeloviti full-stack projekt. Svako predavanje dodaje novi sloj na postojeći rad.
\end{frame}

% ------------------------------------------------------------
\begin{frame}{Završne napomene}
\begin{itemize}
  \item Projekt je \textbf{vaš} - tema je vaš izbor, ali zahtjevi su jasni
  \item Koristite AI za učenje, ne za izbjegavanje razmišljanja
  \item Kod koji ne razumijete je kod koji ne možete obraniti
  \item Pitajte kad zapnete - na predavanjima ili mailom
\end{itemize}

\vspace{0.8em}
\begin{center}
\Large Pitanja?
\end{center}
\end{frame}

\end{document}
